\section{The prime distribution shared by the two analytic routes}
\label{sec:sr-shared-distribution}
The arithmetic-geometric approach controls the smoothed error \(E\).
A spectral approach seeks a trace identity for its prime-power
distribution. The following calculation identifies the common object,
including its normalization, without assuming RH or a trace formula.

\subsection{An exact identity on the logarithmic line}
Set \(H(u)=e^{-u/2}E(e^u)\) for \(u\geq0\), and let \(D_u=d/du\).
For \(u>0\), in distributions,
\begin{equation}\label{eq:sr-prime-distribution}
 (D_u^2-\tfrac14)H
  =\sum_{n\geq2}\frac{\Lambda(n)}{\sqrt n}\delta_{\log n}
          -e^{u/2}\,du .
\end{equation}
This sum is locally finite. In particular, for every
\(\varphi\in C_c^\infty(0,\infty)\),
\begin{equation}\label{eq:sr-prime-test}
 \int_0^\infty H(u)(\varphi''(u)-\tfrac14\varphi(u))\,du
 =\sum_{n\geq2}\frac{\Lambda(n)}{\sqrt n}\varphi(\log n)
       -\int_0^\infty e^{u/2}\varphi(u)\,du.
\end{equation}

Here is a direct proof. Put \(a=\log n\). The contribution of \(n\)
to \(e^{-u/2}T(e^u)\) is
\[
 h_n(u)=\frac{\Lambda(n)}{\sqrt n}
 \bigl(e^{(u-a)/2}-e^{-(u-a)/2}\bigr){\bf1}_{u\geq a}.
\]
It is continuous, is zero before \(a\), and solves
\(h_n''-\tfrac14 h_n=0\) away from \(a\).
Its first derivative jumps at \(a\) by \(\Lambda(n)/\sqrt n\).
Splitting the integral at \(a\) and integrating by parts twice
therefore gives
\[
 \int_0^\infty h_n(u)(\varphi''(u)-\tfrac14\varphi(u))\,du
       =\frac{\Lambda(n)}{\sqrt n}\varphi(a).
\]
Only finitely many \(h_n\) contribute on the support of \(\varphi\).
The remaining term in \(H\) is
\[
 (1+\gamma-u)e^{u/2},\qquad
 (D_u^2-\tfrac14)((1+\gamma-u)e^{u/2})=-e^{u/2}.
\]
This proves \eqref{eq:sr-prime-test} and
\eqref{eq:sr-prime-distribution}. There is no boundary term at zero,
since the test functions vanish in a neighborhood of zero.

The values at zero, useful for a one-sided Laplace transform, are
\[
 H(0)=1+\gamma,\qquad H'(0)=\frac{\gamma-1}{2}.
\]
For \(\Re z>3/2\) the elementary growth bound already proved for
\(E(e^u)\) justifies integration by parts at infinity.
Changing variables gives
\(\mathcal L H(z)=\Phi(z+\tfrac12)\), and
\eqref{eq:sr-prime-distribution} yields
\begin{align}
 (z^2-\tfrac14)\Phi(z+\tfrac12)
  &=-\frac{\zeta'(z+\tfrac12)}{\zeta(z+\tfrac12)}
       -\frac1{z-\tfrac12}\notag\\
  &\qquad +(1+\gamma)z+\frac{\gamma-1}{2}.
 \label{eq:sr-centered-transform}
\end{align}
The boundary terms have been displayed, so they cannot be lost by
formally transforming a differential operator on a half-line.
Equation~\eqref{eq:sr-mellin} gives the same identity algebraically.
The poles associated to nontrivial zeros are located at
\(z=\rho-\tfrac12\). Thus a bound on \(E\) and a valid global spectral
identity address the same zeros, by different additional arguments.
The local prime-power calculation itself proves neither of them.

\subsection{The growth estimate actually needed for spectral sums}
The following global zero count is enough for smooth compactly
supported tests; a universal logarithmic bound in each unit-height
box for arbitrary Mellin transforms is not needed.
\begin{srlemma}\label{lem:sr-global-zero-count}
The nontrivial zeros of \(\zeta\), counted with multiplicity, satisfy
\[
 \#\{\rho:|\rho|\leq R\}=O(R\log(R+2)).
\]
\end{srlemma}
\begin{proof}
From \eqref{eq:sr-reflection},
\[
 \xi(s)=\frac{s(s-1)}2\pi^{-s/2}\Gamma(s/2)\zeta(s)
  =\frac12+\frac{s(s-1)}4J(s),
\]
where
\[
 J(s)=\int_1^\infty(\Theta(t)-1)
           (t^{s/2-1}+t^{(1-s)/2-1})\,dt
\]
is entire and \(\xi(0)=1/2\).
For \(t\geq1\),
\[
 0\leq\Theta(t)-1
 \leq\frac{2e^{-\pi t}}{1-e^{-\pi}}.
\]
Indeed \(n^2\geq n\) bounds the positive-index sum by a geometric
series. For \(|s|\leq R\), \(R\geq2\), both powers in \(J\)
are bounded by \(t^m\), with \(m=\lceil(R+1)/2\rceil\).
Consequently
\[
 |J(s)|\leq C\int_1^\infty e^{-\pi t}t^m\,dt
       \leq C\,\frac{m!}{\pi^{m+1}}
       \leq C(R+2)^{R+2}.
\]
The integral over \([0,\infty)\) equals \(m!/\pi^{m+1}\), by
\(m\) integrations by parts.
Hence
\(\log\max_{|s|\leq R}|\xi(s)|\leq C'(R+2)\log(R+2)\).

For completeness the needed Jensen inequality follows by factoring
the finitely many zeros in a closed disk and applying the mean-value
property to the logarithm of the remaining nonvanishing holomorphic
factor. Its contribution from a zero \(\rho\) in the disk of radius
\(R_1\) is \(\log(R_1/|\rho|)\); this follows from
the mean of \(\log|R_1e^{i\theta}-\rho|\), which is \(\log R_1\).
Choose \(2R<R_1<2R+1\) with no zero on its boundary.
Each zero with \(|\rho|\leq R\) then contributes at least \(\log2\).
The sum of these contributions is bounded by
\(\log\max_{|s|\leq R_1}|\xi(s)|-\log|\xi(0)|\).
This proves the assertion for zeros of \(\xi\).
Every nontrivial zeta zero occurs in \(\xi\) with the same multiplicity,
since the other factors are holomorphic and nonzero there.
\end{proof}

\subsection{What an exact spectral identity on shrinking strips would prove}
\begin{srlemma}[Uniqueness for a vertical discrete spectrum]
\label{lem:sr-spectrum-unique}
Let \(\mathcal Z\) be a closed discrete set in
\(\{z:|\Re z|\leq B\}\).
Give its points nonzero multiplicities \(m_z\) such that
\(\sum_{z\in\mathcal Z}|m_z|(1+|\Im z|)^{-M}<\infty\)
for some \(M\geq0\).
For \(\varphi\in C_c^\infty(0,\infty)\) put
\(F_\varphi(z)=\int_0^\infty\varphi(u)e^{zu}\,du\).
If
\[
 \sum_{z\in\mathcal Z}m_zF_\varphi(z)=0
       \quad\hbox{for every such }\varphi,
\]
then \(\mathcal Z\) is empty.
\end{srlemma}
\begin{proof}
Integration by parts an arbitrary number of times in \(u\)
shows rapid decay of \(F_\varphi(z)\) as \(|\Im z|\to\infty\),
uniformly for \(|\Re z|\leq B\).
There are no boundary terms, because the support is compact in
\((0,\infty)\). Thus
\(a_z=m_zF_\varphi(z)\) is absolutely summable.

For \(T\geq0\), define
\(\varphi_T(u)=\varphi(u-T)\), with zero extension.
It is another permitted test, and
\(F_{\varphi_T}(z)=e^{zT}F_\varphi(z)\).
The hypothesis becomes \(\sum_z a_ze^{zT}=0\).
For \(\Re w>B\), multiply by \(e^{-wT}\) and integrate \(T\)
from zero to infinity. The sum of the absolute integrals is at
most \(\sum_z|a_z|/(\Re w-B)\), so interchange is justified.
We obtain
\[
 \sum_{z\in\mathcal Z}\frac{a_z}{w-z}=0\qquad(\Re w>B).
\]
This series converges normally on every compact set disjoint from
\(\mathcal Z\): only finitely many \(z\) are near that compact set,
and \(|w-z|\geq1\) uniformly for all remaining sufficiently large
\(z\). It therefore defines a meromorphic function, with residue
\(a_z\) at \(z\).
The plane minus a closed discrete set is connected. One can join
two points by a polygonal path and detour around the finitely many
points on its compact image.
The identity theorem makes this meromorphic function zero off
\(\mathcal Z\), so all its residues vanish.

If \(z_0=\sigma+i\tau\) belonged to \(\mathcal Z\), choose a nonzero
nonnegative smooth bump \(b\) supported in \((0,\infty)\) and set
\(\varphi(u)=b(u)e^{-i\tau u}\).
Then \(F_\varphi(z_0)=\int b(u)e^{\sigma u}\,du>0\).
This contradicts \(a_{z_0}=0\), since \(m_{z_0}\ne0\).
\end{proof}

Suppose, as an additional hypothesis, that for every \(\delta>0\)
one had an identity
\begin{equation}\label{eq:sr-strip-trace}
 \mathcal A(\varphi)=
  \sum_{\substack{\rho\\|\Re\rho-1/2|<\delta}}
        m_\rho F_\varphi(\rho-\tfrac12)
       \qquad(\varphi\in C_c^\infty(0,\infty)),
\end{equation}
with exactly the same functional \(\mathcal A\) at every width.
The zero count in Lemma~\ref{lem:sr-global-zero-count} makes these
sums absolutely convergent, by integration by parts as above.
Subtracting the identities at \(0<\delta_1<\delta_2\) gives the
zero distribution for
\[
 \{\,\rho-\tfrac12:
             \delta_1\leq|\Re\rho-\tfrac12|<\delta_2\,\}.
\]
This is closed discrete, because zeros have no finite accumulation
point. Lemma~\ref{lem:sr-spectrum-unique} makes it empty.
Any off-line zero would belong to such a set for some two widths,
so \eqref{eq:sr-strip-trace} would imply RH.

This implication explains the burden of a width-independent trace
claim. Local calculations of the prime and gamma terms do not prove
\eqref{eq:sr-strip-trace}. A construction must prove the spectral
identification, convergence and equality of the global traces,
including any terms introduced by changing the strip.
This paper does not assume \eqref{eq:sr-strip-trace} in its
arithmetic-geometric argument.

\subsection{Two checks on proposed replacements for that burden}
For \(\Re A>0\) put
\[
 g_A(x)=(4\pi A)^{-1/2}\exp(-x^2/(4A)),
 \qquad h(x)=\frac{g_{1+i}(x)-g_{1-i}(x)}{2i}.
\]
Use the square root continuous in the right half-plane and positive
on the positive real axis. Then \(h\) is real-valued.
For every real \(c\), every derivative of \(e^{cx}h(x)\) is a
polynomial times a Gaussian with strictly negative quadratic real
part. It is therefore Schwartz.
Completing the square for positive real \(A\) gives
\(\int_{\mathbb R}g_A(x)e^{zx}\,dx=e^{Az^2}\).
Both sides are holomorphic for \(\Re A>0\), locally uniformly by
Gaussian domination; continuation in \(A\) proves the same formula
there. Thus, for \(\varphi(t)=h(\log t)\),
\[
 \int_0^\infty\varphi(t)t^z\,\frac{dt}{t}
                  =e^{z^2}\sin(z^2).
\]
This entire transform has simple zeros at
\(z=i\sqrt{n\pi}\), \(n\geq1\).
The number in \([-1,1]+i[T,T+1]\), for large \(T\), equals
\[
 \left\lfloor\frac{(T+1)^2}{\pi}\right\rfloor
 -\left\lceil\frac{T^2}{\pi}\right\rceil+1
                  =\frac{2T}{\pi}+O(1).
\]
It cannot satisfy an \(O(\log T)\) unit-box bound.
Rapid vertical decay alone consequently cannot justify that bound
for arbitrary exponential-Schwartz Mellin transforms.
In Jensen's inequality an upper bound on a rapidly decaying value
does not supply the lower bound on that value needed to bound its
negative logarithm.

There is also an elementary projection check.
Let \(\mathcal H\) be a Hilbert space and \(\lambda_0\) a nonzero
continuous linear functional.
If the range of \(P_T\) lies in \(\ker\lambda_0\) for every
\(T\geq T_0\), then \(P_T\) cannot converge strongly to the identity:
otherwise, for every \(x\),
\(\lambda_0(x)=\lim_T\lambda_0(P_Tx)=0\).
More generally intersections of an increasing collection of
evaluation kernels decrease as the collection grows.
Passing to complementary projectors changes the limiting subspace;
it does not prove density in the full space.

These two checks apply to the general zero-count and kernel-projector
claims in Masullo's current draft, Lemma~9.3.2.4 and
Theorem~9.3.2.6 (printed pages 425--426).
The affected statements need replacement before their downstream
spectral argument can supply an input here.
This identifies particular obstructions; it is not a verdict on every
other result in that draft.
The source PDF and the separate, longer audit are listed in
Section~\ref{sec:sr-source-review}.
